% Exploratory preprint. Numeric results are generated by scripts/make_paper.py.
\documentclass[10pt,twocolumn]{article}
\usepackage[margin=0.8in,columnsep=0.27in]{geometry}
\usepackage[T1]{fontenc}
\usepackage{booktabs,graphicx,hyperref,amsmath,times}
\usepackage{xurl}
\hypersetup{colorlinks=true,urlcolor=blue,citecolor=blue,linkcolor=blue}
\input{generated/numbers}
\title{Early Posting Behaviour and Later Deep Engagement in Incel Communities on Reddit: An Exploratory Observational Study}
\author{Raju Sah\\Independent researcher\\\small\href{https://orcid.org/0009-0006-0412-6139}{ORCID: 0009-0006-0412-6139}}
\date{October 24, 2026\\\small Exploratory preprint, version 0.1.0}
\begin{document}
\maketitle

\begin{abstract}
We examine whether early posting behaviour is associated with sustained subsequent participation in three incel-related Reddit communities. The partial archive contains \NKept{} cleaned contributions. Of \NCohort{} first-observed contributors, \NEligA{} have observable primary follow-up labels. Using temporally separated users, an activity-count logistic regression reaches test AUROC \AurocCIAct{}, while a model using posting counts, timing and reply structure reaches \AurocCILgbm{}. The test set contains \NTest{} users, including \PosTest{} positives (\RateTest\%). We report precision-recall performance, F1 and calibration alongside ranking results. Scores, removal status and text are excluded from the corrected analysis because their historical availability or privacy properties could not be established. Earlier test results informed corrections to the cohort and analysis, so these results are exploratory and require independent replication. Community activity is an operational outcome, not a measure of ideology, radicalization or individual dangerousness.
\end{abstract}

\section{Introduction}
Research on the manosphere has documented community growth, overlap and migration across online forums and Reddit \cite{ribeiro2021}. Those aggregate patterns motivate a narrower question: among first-observed contributors to selected communities, how informative is early posting behaviour about later sustained participation?

We study observed engagement rather than ideological change. The outcome requires both a minimum contribution count and activity across several weeks after the feature window. Posting in a target community does not establish agreement with its ideas. Our analyses support retrospective aggregate research and provide no per-user prediction export or operational scoring service.

This preprint presents a corrected, metadata-only analysis, mandatory trivial baselines, temporal evaluation, calibration and sensitivity to the engagement definition. A local analysis plan preceded the original model runs. Subsequent data-validity corrections followed inspection of earlier test results. We therefore describe the final work as exploratory rather than an untouched confirmatory test of that plan.

\section{Related work}
Horta Ribeiro et al.\ \cite{ribeiro2021} characterize the evolution of manosphere communities and migration between them. Their work concerns ecosystem dynamics; our question concerns sustained participation following first observed entry within a restricted archive.

Danescu-Niculescu-Mizil et al.\ \cite{danescu2013} connect early linguistic adaptation and community lifespan. Our corrected analysis instead restricts predictors to posting metadata. It does not test whether language adds information once behaviour is held constant.

Studies of platform interventions, including Reddit bans \cite{chandrasekharan2017}, show why observation periods and community closures matter when interpreting activity. Our design estimates predictive associations and makes no causal claim about bans or migration. Internet research ethics also requires attention to context, identifiability and potential harm across the research lifecycle \cite{aoir2020}.

\section{Data and ethics}\label{sec:ethics}
\paragraph{Sources and snapshot.} Public Reddit submissions and comments were acquired through the Arctic Shift subreddit-level API \cite{arcticshift}. The targets are r/incels, r/Braincels and r/IncelsWithoutHate. Comparison communities and six gateway communities were included in the acquisition plan. No author-history or profile endpoint was queried. Acquisition stopped at a partial snapshot after repeated API errors. Table~\ref{tab:streams} describes its streams.

Ingestion read \NRead{} records and retained \NKept{} after cleaning, with \NAuthorsAll{} distinct hashed authors across all included communities. Authors with deleted identifiers and accounts matching configured bot heuristics are excluded. Contributions from known authors remain activity records when their content is empty or removed. Archive completeness here means completion of the configured retrieval, not proof that every Reddit contribution was captured.

\paragraph{Privacy and minimization.} Authors are replaced with salted SHA-256 hashes. Hashing is pseudonymization and does not guarantee anonymity. The original ingestion retained free text, which can contain identifiers despite author-field hashing. During release preparation, the text fields in the local interim files were cleared, and ingestion was changed to discard text before persistence. The corrected models use no text, score or text-presence predictor. Raw archives remain confined to local raw-data storage. The release excludes raw data, hashes, record identifiers, individual feature rows, predictions, quotations and individual explanations. Published figures summarize groups or feature-level averages.

\paragraph{Ethics review.} No formal ethics review was sought for this study. This statement does not represent an institutional approval or exemption determination. Public accessibility alone does not remove ethical concerns. Following the risk considerations discussed by AoIR \cite{aoir2020}, we limit publication to aggregate findings and avoid quotations, named-user investigation and person-level outputs. The absence of independent ethics review remains a limitation.

\section{Methods}
\paragraph{Entry and feature window.} Entry $t_0$ is the earliest cleaned contribution observed in a target community, after deduplication by record kind and identifier. It is not necessarily the user's true first participation. The Task A cutoff is the earlier of the fifth target contribution and 14 days after $t_0$. Feature records satisfy $t_0\leq t\leq\mathrm{cutoff}$.

\paragraph{Outcome and eligibility.} The primary outcome requires at least 20 target contributions in at least four distinct seven-day bins during $(\mathrm{cutoff},\mathrm{cutoff}+90\ \mathrm{days}]$. Bins are anchored immediately after the cutoff. Comments and submissions have equal weight. Windows extending past the configured observed coverage of a target considered present at entry receive null labels. Coverage bounds summarize observed records and do not prove the absence of internal archive gaps or verify a community's actual lifetime.

\paragraph{Corrected predictors.} Predictors are contribution count, comment share, active days, inter-event gap statistics, UTC hour entropy, weekend share, distinct observed thread and target-community counts, reply-to-comment share, time to the fifth contribution, and whether that contribution was reached. The activity baseline uses contribution count alone. Creation-time filtering is tested explicitly. Scores, deletion or removal status, sentiment, text lengths, TF-IDF and embeddings are excluded because record creation dates do not establish when mutable fields attained their archived values. The corrected design addresses that identified availability problem; it does not establish that all sources of bias or archive error are absent.

\paragraph{Temporal splits.} The corrected analysis retains the round-two cohort and split. Entrants after 2018-09-19 are excluded. Split boundaries are computed among primary-label-observable users with nominal fractions 0.50/0.25/0.25. Users sharing a UTC entry date remain together. A 90-day follow-up embargo excludes earlier-split users whose label windows reach the next split's start. Realized split sizes differ from the nominal fractions after embargo. Table~\ref{tab:balance} reports dates and class balance. These changes followed earlier test inspection and are disclosed as D3 and D4.

\paragraph{Models and selection.} We compare a majority classifier, activity-count logistic regression, metadata logistic regression and LightGBM \cite{ke2017}. The majority classifier uses training prevalence as a constant probability and a fixed 0.5 decision threshold, so it predicts the majority training class. Logistic regressions standardize predictors using training data and use balanced class weights. LightGBM uses a training-derived positive-class weight. Users are never resampled before splitting. The original validation tuning grids and seed are retained in \texttt{configs/release.yaml}. Hyperparameters maximize validation average precision. Nontrivial-model thresholds maximize validation F1 over a fixed quantile grid. Each corrected model is scored once per declared analysis; existing result files are reused on a rerun only when input, code and configuration provenance match.

\paragraph{Evaluation.} We report AUROC, average precision as AUPRC, F1 at the specified threshold, expected calibration error (ECE, ten equal-width bins), Brier score and reliability curves. Precision-recall metrics are useful alongside ROC metrics under class imbalance \cite{saito2015}. Brier score measures overall probabilistic performance rather than calibration alone \cite{sklearncalibration}. Percentile confidence intervals use 1,000 user-level bootstrap resamples with a fixed seed, omitting resamples containing only one class. Intervals condition on the fitted models and observed cohort; they do not incorporate uncertainty from archive gaps, model selection or post-hoc design choices. Reuse of an earlier-inspected test population means this is not an independent confirmation.

\section{Results}
\paragraph{Cohort.} The cohort contains \NCohort{} users, of whom \NEligA{} are primary-label-observable and \NDeepA{} meet the primary outcome (\PctDeepA\%). After eligibility and embargo, training contains \NTrain{} users (\PosTrain{} positives, \RateTrain\%), validation \NVal{} (\PosVal{}, \RateVal\%) and test \NTest{} (\PosTest{}, \RateTest\%).

\begin{table*}[t]\centering\small
\input{generated/tab_balance}
\caption{Primary-label class balance and UTC entry-date ranges after embargo.}\label{tab:balance}
\end{table*}

\begin{table*}[t]\centering\small\setlength{\tabcolsep}{5pt}
\input{generated/tab_release}
\caption{Corrected metadata-only Task A results. Brackets give conditional 95\% bootstrap intervals. ECE and Brier score are better when lower. The constant baseline retains training-prevalence probabilities and uses a fixed majority decision.}\label{tab:main}
\end{table*}

\paragraph{Primary analysis.} The activity-only baseline achieves AUROC \AurocAct{} and AUPRC \AuprcAct{}. Metadata LightGBM achieves AUROC \AurocLgbm{} and AUPRC \AuprcLgbm{} (Table~\ref{tab:main}). These are descriptive comparisons across fixed model families, not a formal paired test of improvement. Class weighting changes the scale of probabilities; calibration must be assessed separately from ranking. All three class-weighted models have larger Brier loss than the constant baseline in this test population. Their scores should not be interpreted as calibrated individual risks. No post-hoc calibrator is fitted or selected using the test set.

\begin{figure}[t]\centering
\includegraphics[width=\columnwidth]{figures/release_reliability}
\caption{Corrected primary models' test reliability curves. Each plotted point summarizes a probability bin containing at least ten users.}\label{fig:cal}
\end{figure}

\paragraph{Validation interpretation.} Figure~\ref{fig:importance} reports mean changes in validation AUPRC when each metadata predictor is permuted. Correlated predictors can share information, so permutation importance does not identify causal mechanisms or establish the absence of useful content information.

\begin{figure}[t]\centering
\includegraphics[width=\columnwidth]{figures/release_importance}
\caption{Validation permutation importance for the selected primary metadata LightGBM. Bars average five permutations. Whiskers show their standard deviation, not confidence intervals.}\label{fig:importance}
\end{figure}

\paragraph{Sensitivity.} The corrected release covers \NRobust{} analyses: nine 90-day threshold combinations, a comments-only outcome and exclusion of near-start entrants. Metadata LightGBM AUROC ranges from \RobAurocLow{} to \RobAurocHigh{} (Table~\ref{tab:rob}). Each scenario retains the original primary cohort split and uses validation selection. These analyses probe operational definitions within the same archive, rather than independent replications. The original nine 180-day variants remain unestimable under the available coverage.

\begin{table*}[t]\centering\small\setlength{\tabcolsep}{4pt}
\input{generated/tab_release_robust}
\caption{Corrected metadata LightGBM sensitivity analyses. All follow-up windows are 90 days.}\label{tab:rob}
\end{table*}

\paragraph{Unresolved planned analyses.} Task B, using pre-entry gateway activity, has only \NTestB{} historical test users and \PosTestB{} positives, and is too limited for a substantive conclusion here. The final-cohort vocabulary ablation has no eligible contrast texts. Its superseded round-one result is exploratory historical material and supports no headline conclusion about community-specific vocabulary. Historical text, embedding and MGTOW results use different predictor availability assumptions or populations. They are not evidence for the corrected analysis, which makes no comparison between text and behaviour.

\section{Limitations}
The archive is incomplete, and coverage limits do not resolve internal gaps or differential missingness. First-observed entry is left-censored, particularly near the effective r/incels start in July 2016. The near-start sensitivity removes a configured 30-day group but cannot restore unobserved history. Bot filtering is heuristic.

The final test set consists of 2018 entrants. Community composition and historical conditions differ across splits, so temporal evaluation combines generalization across time with community and population shifts. The observed association is specific to selected subreddits, this snapshot and the chosen engagement thresholds. Activity outside those targets is unobserved. Records with removed content can still contribute activity counts when an author is known.

Early contribution volume and the outcome both measure activity. Their association may reflect persistence in posting rather than a distinct ideological pathway. No inference about radicalization, extremist beliefs, intent, dangerousness or causal influence follows from this label.

The cohort restriction, split changes and corrected feature set followed inspection of earlier test results. Freezing this release configuration prevents further within-run test tuning, but does not remove researcher degrees of freedom or recover an untouched holdout. Independent data and prospective choices are needed for confirmation. No ethics committee reviewed the study.

\section{Conclusion}
Early posting metadata is associated with later sustained target-community participation in this partial Reddit archive. The corrected analysis excludes mutable archive fields and text, reports trivial baselines and calibration, and documents changes prompted by validity and privacy concerns. The findings are exploratory aggregate associations. They do not establish ideological change or justify acting on individuals.

\section*{Declarations}
\paragraph{Author and funding.} Raju Sah is the sole author and an independent researcher. This research received no external funding. The author declares no competing interests.
\paragraph{AI assistance.} OpenAI Codex and Anthropic Claude assisted with code development, analysis support, manuscript preparation and review. Raju Sah retains responsibility for the research decisions and the final manuscript. This disclosure does not substitute for author verification of the code, evidence and claims.
\paragraph{Availability and licenses.} Supporting code and synthetic tests accompany the preprint under the MIT license. The manuscript is licensed under Creative Commons Attribution 4.0 International. Raw Reddit archives, derived user-level tables, hashes, text, embeddings and predictions are excluded. Aggregate tables and figures appear in this manuscript. Reproducing numerical results requires separately authorized access to the same partial source snapshot; code and tests alone cannot reconstruct unavailable records. The release instructions distinguish regenerating the local paper from running a new analysis.

\begin{thebibliography}{9}
\bibitem{ribeiro2021} Manoel Horta Ribeiro, Jeremy Blackburn, Barry Bradlyn, Emiliano De Cristofaro, Gianluca Stringhini, Summer Long, Stephanie Greenberg, and Savvas Zannettou. 2021. The Evolution of the Manosphere across the Web. \emph{Proceedings of ICWSM}, 15(1):196--207. \url{https://doi.org/10.1609/icwsm.v15i1.18053}.
\bibitem{danescu2013} Cristian Danescu-Niculescu-Mizil, Robert West, Dan Jurafsky, Jure Leskovec, and Christopher Potts. 2013. No Country for Old Members: User Lifecycle and Linguistic Change in Online Communities. \emph{Proceedings of WWW}, 307--318. \url{https://doi.org/10.1145/2488388.2488416}.
\bibitem{chandrasekharan2017} Eshwar Chandrasekharan, Umashanthi Pavalanathan, Anirudh Srinivasan, Adam Glynn, Jacob Eisenstein, and Eric Gilbert. 2017. You Can't Stay Here: The Efficacy of Reddit's 2015 Ban Examined Through Hate Speech. \emph{Proceedings of the ACM on Human-Computer Interaction}, 1(CSCW), Article 31. \url{https://doi.org/10.1145/3134666}.
\bibitem{aoir2020} aline shakti franzke, Anja Bechmann, Michael Zimmer, Charles M. Ess, and the Association of Internet Researchers. 2020. \emph{Internet Research: Ethical Guidelines 3.0}. \url{https://aoir.org/reports/ethics3.pdf}.
\bibitem{arcticshift} Arctic Shift contributors. \emph{Project Arctic Shift}. Project documentation, accessed October 24, 2026. \url{https://github.com/ArthurHeitmann/arctic_shift}.
\bibitem{ke2017} Guolin Ke, Qi Meng, Thomas Finley, Taifeng Wang, Wei Chen, Weidong Ma, Qiwei Ye, and Tie-Yan Liu. 2017. LightGBM: A Highly Efficient Gradient Boosting Decision Tree. \emph{Advances in Neural Information Processing Systems}, 30. \url{https://papers.nips.cc/paper_files/paper/2017/hash/6449f44a102fde848669bdd9eb6b76fa-Abstract.html}.
\bibitem{saito2015} Takaya Saito and Marc Rehmsmeier. 2015. The Precision-Recall Plot Is More Informative than the ROC Plot When Evaluating Binary Classifiers on Imbalanced Datasets. \emph{PLOS ONE}, 10(3):e0118432. \url{https://doi.org/10.1371/journal.pone.0118432}.
\bibitem{sklearncalibration} scikit-learn developers. \emph{Probability calibration}. Version 1.9.1 documentation, accessed October 24, 2026. \url{https://scikit-learn.org/stable/modules/calibration.html}.
\end{thebibliography}

\appendix
\section{Data snapshot}
\begin{table*}[t]\centering\small
\input{generated/tab_streams}
\caption{Local acquisition snapshot. Completion describes retrieval status, not verified completeness of Reddit history.}\label{tab:streams}
\end{table*}
\section{Analysis changes}
\texttt{DEVIATIONS.md} records D1 (partial acquisition and coverage overrides), D2 (separate MGTOW split), D3 (entry cutoff), D4 (eligible-only split and 90-day embargo), D5 (near-start anchor), and D6 (corrected metadata-only release analysis, majority baseline and privacy minimization). The earlier amendment moved MGTOW from the target list to the comparison list. The final results should be interpreted alongside that history.
\end{document}
